\section{Universal tensor laws for the saturating states}
\label{sec:universal-tensor}

The single-copy saturation normal form has consequences even when the
second tensor factor is an arbitrary mixed state. Throughout, tensor
products are regrouped as $(A_1A_2):(B_1B_2)$, rather than as the two
original bipartite systems. For a coefficient matrix $C$, write
$|C\rangle=\sum_{ij}C_{ij}|i,j\rangle$; its Schmidt rank is $\operatorname{rank}C$.
Ensembles below use subnormalized vectors, so that
$\rho=\sum_j|C_j\rangle\langle C_j|$ without separate probabilities.

\paragraph{Idea.} By the saturation normal form, every saturator contains a
fixed coherent vector $z$ that is orthogonal to the residual. Projecting any
ensemble of $\rho\otimes\sigma$ onto $z$ in the first factor returns an
ensemble of $\sigma$, and the coefficient matrices of the ensemble vectors
have a block-triangular form whose diagonal blocks are multiples of those of
the projected vectors. Ranks and determinants therefore multiply.

\begin{theorem}[Universal Schmidt-number multiplication]
\label{thm:universal-schmidt}
Let $\rho$ be any state saturating the sharp marginal-spectrum fully
entangled fraction bound, and put
$r=\min(\operatorname{rank}\rho_A,\operatorname{rank}\rho_B)$.
For every finite-dimensional bipartite state $\sigma$, with arbitrary
rectangular local dimensions,
\[
 \boxed{\quad \operatorname{SN}(\rho\otimes\sigma)
                      =r\,\operatorname{SN}(\sigma).\quad}
\]
In particular $\operatorname{SN}(\rho^{\otimes n})=r^n$ for every
positive integer $n$.
\end{theorem}
\begin{proof}
Use the established saturation normal form
\[
 \rho=|z\rangle\langle z|+R,\qquad
 z=\sum_i\sqrt{m_i}|ii\rangle,\qquad
 \operatorname{supp}R\subseteq\mathbb C^I\otimes\mathbb C^J,
 \qquad I\cap J=\varnothing.
\]
Here precisely $r$ entries of $m$ are positive and $L=\|z\|^2>0$.
Every vector in $\operatorname{supp}(\rho\otimes\sigma)$ has a unique
coefficient along the first factor $z$: it can be written
\[
       |v\rangle=|z\rangle\otimes|v_0\rangle+|v_\perp\rangle,
       \qquad v_\perp\in\operatorname{supp}R\otimes\operatorname{supp}\sigma.
\]
The coefficient vector is
$v_0=(\langle z|/L\otimes I)v$.
This is an algebraic ensemble-extraction map; it is not asserted to be a
local physical operation.

If $C_0$ is the coefficient matrix of $v_0$, the coefficient matrix $T$
of $v$ has block form
\[
       T=\operatorname{diag}(\sqrt m)\otimes C_0
                        +\sum_{i\in I,j\in J}E_{ij}\otimes C_{ij}.
\]
Order the positive-cap indices with $I$ first and $J$ last. The principal
first-factor $r\times r$ block array is upper triangular, with diagonal
blocks $\sqrt{m_i}C_0$. If $C_0$ has rank $k$, select a nonzero $k\times k$
minor and repeat those row and column selections in each diagonal block.
The resulting $rk\times rk$ minor of $T$ has determinant
\[
        \left(\prod_{i:m_i>0}m_i^{k/2}\right)
                               \bigl(\det C_0[\mathcal I,\mathcal J]\bigr)^r
                           \ne0.
\]
Consequently $\operatorname{rank}T\ge r\operatorname{rank}C_0$, including
rectangular $C_0$.

Now take any subnormalized ensemble
$\rho\otimes\sigma=\sum_\ell|v_\ell\rangle\langle v_\ell|$.
Every ensemble vector lies in the support of the state. Applying the
coefficient-extraction map on both sides gives exactly
\[
       \sum_\ell|v_{0,\ell}\rangle\langle v_{0,\ell}|=\sigma,
\]
because $Rz=0$ and $\langle z|\rho|z\rangle=L^2$.
By the definition of Schmidt number, at least one nonzero $v_{0,\ell}$
has Schmidt rank at least $\operatorname{SN}(\sigma)$. The preceding
minor proves the required lower bound for every ensemble of
$\rho\otimes\sigma$.
For the upper bound, take ensembles of $\rho$ and $\sigma$ with maximal
Schmidt ranks $r$ and $\operatorname{SN}(\sigma)$, respectively, and use
their product ensemble. Pure-state Schmidt ranks multiply. This proves
equality and induction gives the tensor-power formula.
\end{proof}

There is also a short derivation from the preceding quasidistillation
result: local filters approach $|\Phi_r\rangle\langle\Phi_r|$ on the first
factor, Schmidt-number sublevel sets are closed, and
$\operatorname{SN}(\Phi_r\otimes\sigma)=r\operatorname{SN}(\sigma)$.
The last identity follows by the same repeated-diagonal-block argument.
Thus the Schmidt-number theorem is a strong consequence of the already
established normal form and quasidistillation; the determinant factorization
of both $G$ roofs below supplies separate additional information.

\begin{theorem}[Universal factorization of both $G$ roofs]
\label{thm:universal-g}
Suppose, in addition, that $\rho$ is on $\mathbb C^d\otimes\mathbb C^d$
and every $m_i>0$. For every state $\sigma$ on
$\mathbb C^e\otimes\mathbb C^e$,
\[
 \boxed{\begin{aligned}
 G_{de}(\rho\otimes\sigma)&=G_d(\rho)G_e(\sigma),\\
 G_{de}^{\rm a}(\rho\otimes\sigma)&=G_d(\rho)G_e^{\rm a}(\sigma),
 \end{aligned}}
 \qquad
 G_d(\rho)=G_d^{\rm a}(\rho)
              =d\left(\prod_i m_i\right)^{1/d}.
\]
There is no flat-roof, rank, marginal, or saturation assumption on
$\sigma$. Here the homogeneous pure-state normalization is
$G_e(C)=e|\det C|^{2/e}$; the convex roof is the infimum of the ensemble
sum and the assistance roof is its supremum.
\end{theorem}
\begin{proof}
Use the coefficient matrix $T$ from the preceding proof. Its $d$ diagonal
blocks have size $e$, so block triangularity gives
\[
       \det T=\left(\prod_i m_i^{e/2}\right)(\det C_0)^d.
\]
This identity also holds when $C_0$ is singular. Therefore, with
$g=d(\prod_i m_i)^{1/d}$,
\[
                   G_{de}(T)=g\,G_e(C_0).
\]
Every ensemble of $\rho\otimes\sigma$ induces an actual ensemble of
$\sigma$ by the exactly normalized projection already proved. It follows
that every ensemble sum lies between $gG_e(\sigma)$ and
$gG_e^{\rm a}(\sigma)$.
Conversely, choose any ensemble of $\rho$ and an optimizing ensemble of
$\sigma$. Every ensemble of $\rho$ has $G$ sum $g$. The determinant
identity for a Kronecker product gives
$G_{de}(C\otimes D)=G_d(C)G_e(D)$, so the product ensemble attains the
corresponding lower or upper endpoint. Equivalently one may use sequences
of ensembles approaching the infimum or supremum; no attainment lemma is
needed for the equalities.
\end{proof}

\subsection{The actual structural hypothesis is a common triangular flag}

The preceding factorization is more general than spectral saturation.
Let $D=\operatorname{diag}(d_1,\ldots,d_d)$ with every $d_i\ne0$, and let
$\mathcal T$ be a linear subspace of strictly upper triangular matrices.
Suppose the coefficient-matrix support of a state $\tau$ is contained in
$\mathbb CD+\mathcal T$ and is not contained in $\mathcal T$.
For $C=aD+N$ define $\ell(C)=a$ and put
\[
       \eta=\sum_j|\ell(C_j)|^2>0
       \quad\text{for any ensemble }\tau=\sum_j|C_j\rangle\langle C_j|.
\]
This number is ensemble independent: it is the expectation of the positive
rank-one operator representing the linear functional $\ell$.
The same arguments give, for arbitrary bipartite $\sigma$,
\[
 \operatorname{SN}(\tau\otimes\sigma)=d\operatorname{SN}(\sigma),
\]
and for square $\sigma$ the two $G$ factorization identities, with
\[
       G_d(\tau)=G_d^{\rm a}(\tau)=d|\det D|^{2/d}\eta.
\]
Indeed the induced subnormalized vectors now sum to $\eta\sigma$;
dividing them by $\sqrt\eta$ gives an ensemble of $\sigma$.
The tensor coefficient matrix still has diagonal blocks $d_iC_0$.
This version permits cross coherences between the distinguished vector
$D$ and the triangular residual subspace. It is also invariant under
invertible local filtering: Schmidt numbers are unchanged, while the
normalized $G$ roofs acquire the common factor
$|\det A\det B|^{2/d}/p$, where $p$ is the filter success probability.
On a tensor product the same factor occurs, so it cancels from each
factorization identity.

\subsection{Exact operational consequences}

For a saturator, the zero-error cost of preparing $n$ copies by LOCC from
a maximally entangled resource is exactly
\[
       E_c^{(0)}(\rho^{\otimes n})=n\log_2 r,
       \qquad E_c^{(0,\infty)}(\rho)=\log_2 r.
\]
The operational identity $E_c^{(0)}(\omega)=\log_2\operatorname{SN}(\omega)$
is established prior work; see Yue and Chitambar,
\href{https://arxiv.org/abs/1808.10516}{arXiv:1808.10516}.
The value here follows by substituting the exact tensor theorem.
This is zero-error preparation cost and makes no equality claim for the
usual vanishing-error asymptotic entanglement cost.

The same theorem rules out erasing a strict Schmidt-number obstruction
with a saturator as an exactly returned, uncorrelated catalyst. If an LOCC
success branch maps $\rho\otimes\sigma$ to $\rho\otimes\tau$ with nonzero
probability, then
$\operatorname{SN}(\tau)\le\operatorname{SN}(\sigma)$.
For a full-positive square saturator the $G$ factorization also gives the
necessary success-probability bound
$pG_e(\tau)\le G_e(\sigma)$, using the standard average monotonicity of
$G$-concurrence. These statements require exact product return of the
catalyst; they do not cover correlated or approximately returned catalysts.
The $G$-concurrence monotonicity and determinant scaling used here are
established properties, reviewed in Gour,
\href{https://arxiv.org/abs/quant-ph/0506229}{arXiv:quant-ph/0506229},
Eqs.~(4)--(5) and the proof of Theorem~1.

Tensor closure of the \emph{saturation condition} itself is not required
and is false. For example set
\[
 z=(3|00\rangle+|11\rangle)/4,
 \qquad \rho_*=|z\rangle\langle z|+\frac38|01\rangle\langle01|.
\]
Its decreasing marginal spectra are
$\alpha=(15,1)/16$, $\beta=(9,7)/16$, and it is a spectral saturator with
$f_2=1/2$ and $G_2=3/8$. But
\[
 \lambda_{\max}(\rho_*^{\otimes2})=\frac{25}{64}
    <\frac7{16}
    =\sum_i\min\{(\alpha^{\otimes2})_i^\downarrow,
                         (\beta^{\otimes2})_i^\downarrow\}.
\]
The earlier equality classification therefore excludes tensor-square
saturation, whereas the present theorem gives Schmidt number four and
$G_4=G_4^{\rm a}=9/64$ exactly.

For an example with different lower and upper roofs in the second factor,
take the two-qubit isotropic state
\[
 \sigma=\frac34|\Phi_2\rangle\langle\Phi_2|
            +\frac1{12}(I-|\Phi_2\rangle\langle\Phi_2|).
\]
Then $G_2(\sigma)=1/2$ and $G_2^{\rm a}(\sigma)=1$, giving
\[
       G_4(\rho_*\otimes\sigma)=\frac3{16},\qquad
       G_4^{\rm a}(\rho_*\otimes\sigma)=\frac38.
\]
The lower value for $\sigma$ follows without a numerical roof search.
For any normalized two-qubit pure state,
$\langle\Phi_2|\psi\rangle\langle\psi|\Phi_2\rangle\le(1+G_2(\psi))/2$,
so every ensemble of $\sigma$ has mean concurrence at least $1/2$.
An attaining six-vector ensemble has coefficient matrices
\[
 C_{j,\pm}=\frac{\sqrt3\,I\pm P_j}{2\sqrt2},
       \qquad j=1,2,3,
\]
with equal weights $1/6$, where the $P_j$ are Pauli matrices.
Each is normalized and has concurrence $1/2$; their mixture is $\sigma$.
The Bell eigenensemble attains assistance one, the universal upper bound.

\section{Flat \texorpdfstring{$G$}{G} roofs can collapse completely under tensoring}
\label{sec:flat-roof-counterexample}

The common triangular flag cannot be replaced by the statement that
the lower and upper $G$ roofs agree. The following explicit counterexample
also shows why a linear space consisting entirely of nilpotent matrices
is not enough for the tensor proof.
The underlying nilpotent pencil is classical: the matrices below are the
transposes of those in Christensen,
\href{https://arxiv.org/abs/funct-an/9702017}{arXiv:funct-an/9702017},
Example~3.1, which credits Wielandt and Motzkin--Taussky.
The quantum state, ensemble conversion and entanglement conclusions are
derived explicitly here; the matrix pencil itself is not claimed as new.

\begin{theorem}[A three-level flat-roof counterexample]
\label{thm:flat-roof-tensor-failure}
Define
\[
 A=\begin{pmatrix}0&1&0\\0&0&1\\0&0&0\end{pmatrix},\qquad
 B=\begin{pmatrix}0&0&0\\1&0&0\\0&-1&0\end{pmatrix},\qquad
 \omega=\frac{|I_3\rangle\langle I_3|
                 +|A\rangle\langle A|+|B\rangle\langle B|}{7}.
\]
This is a rank-three state on $\mathbb C^3\otimes\mathbb C^3$, with
\[
 G_3(\omega)=G_3^{\rm a}(\omega)=\frac37,
       \qquad \operatorname{SN}(\omega)=3.
\]
Nevertheless,
\[
 \boxed{\quad G_9(\omega\otimes\omega)=0,\qquad
          \operatorname{SN}(\omega\otimes\omega)\le7<9,\qquad
          G_9^{\rm a}(\omega\otimes\omega)\ge\frac9{49}>0.\quad}
\]
The bound seven is an explicit ensemble upper bound; equality of the
tensor-square Schmidt number with seven is not asserted.
\end{theorem}
The nilpotent plane $\operatorname{span}\{A,B\}$ has no common triangular
flag: $\ker A=\mathbb Ce_0$, $\ker B=\mathbb Ce_2$, so their common
kernel is zero. Simultaneously strictly upper triangular matrices would
annihilate the same nonzero first basis vector.
\begin{proof}
The three coefficient matrices $I_3,A,B$ are Hilbert--Schmidt orthogonal
with squared norms $3,2,2$. Positivity and trace one are therefore
immediate, and the eigenvalues are $3/7,2/7,2/7$.
Direct multiplication gives
\[
        (aA+bB)^3=0,\qquad
        \det(cI_3+aA+bB)=c^3
           \quad(a,b,c\in\mathbb C).
\]
In any subnormalized ensemble of $\omega$, write the coefficient
matrices as $C_j=c_jI_3+a_jA+b_jB$. Taking the expectation on
$|I_3\rangle/\sqrt3$ gives $\sum_j|c_j|^2=1/7$. Hence
\[
         \sum_jG_3(C_j)=3\sum_j|c_j|^2=3/7
\]
for \emph{every} ensemble. This proves both roof values. Every ensemble
has a vector with $c_j\ne0$, hence of Schmidt rank three, proving the
Schmidt-number assertion.

The tensor product initially has the nine-term subnormalized ensemble
$|C\otimes D\rangle/7$, for $C,D\in\{I_3,A,B\}$.
Let $\zeta=e^{2\pi i/3}$. Replace the three coefficient matrices
$I_3\otimes I_3,A\otimes B,B\otimes A$ by
\[
          V_j=\frac{I_9+\zeta^jA\otimes B
                           +\zeta^{2j}B\otimes A}{\sqrt3},
                   \qquad j=0,1,2.
\]
The $3\times3$ Fourier matrix is unitary, so this substitution preserves
the sum of the three vector projectors exactly. The remaining six terms
are left unchanged.
The determinant identity
\[
 \det(I_9+xA\otimes B+yB\otimes A)
                       =(xy-1)^2(2xy+1)
\]
shows that every $V_j$ is singular. More precisely, for $xy=1$ the
matrix has the two independent kernel vectors
\[
    -x|01\rangle+|10\rangle,
                 \qquad x|12\rangle+|21\rangle.
\]
Its principal minor indexed, in zero-based lexicographic order, by
$\{0,1,2,4,5,6,8\}$ has determinant three. Thus its rank is exactly
seven for $x\ne0$, and this applies to all three Fourier vectors.
The six remaining matrices are
$I\otimes A,I\otimes B,A\otimes I,B\otimes I,A\otimes A,B\otimes B$;
their ranks are $6,6,6,6,4,4$.

After normalization, the three rank-seven ensemble vectors each have
probability $17/147$, the four rank-six vectors each have probability
$6/49$, and the two rank-four vectors each have probability $4/49$.
These probabilities sum to one. This is an exact ensemble of
$\omega\otimes\omega$ in which every coefficient determinant vanishes,
so the convex-roof $G_9$ is zero and the Schmidt number is at most seven.
The original product eigenensemble has $G_9$ sum $(3/7)^2$, proving the
positive assistance lower bound.
\end{proof}

For additional checks, both actual marginals of $\omega$ are
$\operatorname{diag}(2,3,2)/7$, its fully entangled fraction is $3/7$,
and its partial-transpose eigenvalues are
\[
       \frac17\ [3\text{ times}],\quad
       \frac27\ [2\text{ times}],\quad
       0\ [2\text{ times}],\quad
       \frac{\sqrt3}{7},\quad-\frac{\sqrt3}{7}.
\]
The FEF value follows because $|I_3\rangle/\sqrt3$ is a maximally
entangled top eigenvector. The state is entangled and NPT, but it is
not a spectral saturator: its equal marginal spectra have sharp FEF
cap $(11+4\sqrt6)/21>3/7$. Equivalently their maximal-eigenvalue cap
is one, which only a pure coupling can attain.

\subsection{A whole family with strict failure of tensor multiplicativity}

The counterexample is not isolated. With the same $A,B$, define
\[
 \omega_p=\frac p3|I_3\rangle\langle I_3|
       +\frac{1-p}{4}\bigl(|A\rangle\langle A|+|B\rangle\langle B|\bigr),
                  \qquad 0<p\le1.
\]
Every member has
\[
       G_3(\omega_p)=G_3^{\rm a}(\omega_p)=p,
       \qquad \operatorname{SN}(\omega_p)=3.
\]
This follows from the same determinant identity and the coefficient
expectation $\sum_j|c_j|^2=p/3$.
For their tensor squares one has the explicit rigorous bounds
\begin{equation}\label{eq:flat-family-bound}
 \boxed{\quad
 G_9(\omega_p^{\otimes2})
       \le\max\left\{0,p^2-\frac9{16}(1-p)^2\right\},
       \qquad G_9^{\rm a}(\omega_p^{\otimes2})\ge p^2.
 \quad}
\end{equation}
In particular the lower roof is exactly zero for $0<p\le3/7$,
and $G_9(\omega_p^{\otimes2})<p^2$ for every $0<p<1$.
For $p>3/7$ the upper bound in \eqref{eq:flat-family-bound} is not claimed
to be the exact convex-roof value.

To prove this, put
\[
 T=\frac{|A\rangle\langle A|+|B\rangle\langle B|}{4},
 \qquad P_3=|\Phi_3\rangle\langle\Phi_3|.
\]
The state $T$ has an explicit ensemble of Schmidt rank two, so
$\operatorname{SN}(T)\le2$. For $0<p\le3/7$,
\[
       \omega_p=\frac{7p}{3}\omega
                      +\left(1-\frac{7p}{3}\right)T.
\]
Its tensor square is a convex combination of
$\omega\otimes\omega$, $\omega\otimes T$, $T\otimes\omega$ and
$T\otimes T$, with Schmidt-number upper bounds $7,6,6,4$.
Thus $\operatorname{SN}(\omega_p^{\otimes2})\le7$, proving the exact
zero lower roof on that interval.
For $3/7\le p\le1$, put $q=(7p-3)/4\in[0,1]$. Then
\[
                   \omega_p=qP_3+(1-q)\omega.
\]
Convexity of the lower roof, the zero tensor-square value for $\omega$,
and the pure-factor tensor identity give
\[
 \begin{aligned}
 G_9(\omega_p^{\otimes2})
  &\le q^2+\frac67q(1-q)\\
  &=\frac{(7p-3)(p+3)}{16}
    =p^2-\frac9{16}(1-p)^2.
 \end{aligned}
\]
The product ensemble supplies the assistance bound $p^2$. At $p=1$ the
state is a pure Bell state and equality with one is recovered. The
strict failure for every nonpure $\omega_p$ is therefore fully explicit.

\subsection{A quantitative obstruction to full-dimensional distillation}

For every one-copy stochastic LOCC output $\tau$ of $\omega$ on a
$3\times3$ output space,
\[
                         f_3(\tau)\le\frac{\sqrt7}{3}<1.
\]
Indeed, perform the same success branch independently on both copies.
It succeeds with positive probability and gives $\tau\otimes\tau$.
Schmidt number cannot increase even on a success branch, so
$\operatorname{SN}(\tau\otimes\tau)\le7$.
Choose a maximizing Bell vector for $\tau$; its tensor square is a
nine-dimensional Bell vector. The standard Schmidt-number overlap bound
$f_D(\nu)\le\operatorname{SN}(\nu)/D$ therefore gives
$f_3(\tau)^2\le7/9$.
This is a rigorous upper bound, not a claim that $\sqrt7/3$ is attainable.
It excludes quasidistillation to a three-dimensional Bell state despite
$G_3=G_3^{\rm a}>0$.
The same bound holds for every $\omega_p$ with $0<p\le3/7$, because the
preceding argument gives the same tensor-square Schmidt-number bound.

In fact \emph{every} nonpure member of the family fails to be
quasidistillable to $|\Phi_3\rangle$. For $0<p<1$ define
\[
 c=\min\left\{\frac{7p}{3},\frac{7(1-p)}4\right\},\qquad
 C=\max\left\{\frac{7p}{3},\frac{7(1-p)}4\right\}.
\]
The common orthogonal eigenvectors give $c\omega\le\omega_p\le C\omega$.
For any stochastic LOCC success map $\Lambda$ with
$q_p=\operatorname{Tr}\Lambda(\omega_p)>0$, positivity gives
$q_0=\operatorname{Tr}\Lambda(\omega)>0$ and $q_p\le Cq_0$.
The normalized outputs consequently satisfy
\[
 \tau_p=\frac{\Lambda(\omega_p)}{q_p}
   \ge\frac cC\frac{\Lambda(\omega)}{q_0}
   =\frac cC\tau_0.
\]
Thus $\tau_p$ is a convex combination with weight $c/C$ on $\tau_0$.
Convexity of the FEF and $f_3(\tau_0)\le\sqrt7/3$ imply the explicit
piecewise bound
\[
 f_3(\tau_p)\le
 \begin{cases}
  \sqrt7/3,&0<p\le3/7,\\[2pt]
  1-\dfrac{3(1-p)}{4p}\left(1-\dfrac{\sqrt7}{3}\right),
       &3/7\le p<1.
 \end{cases}
\]
Both expressions agree at $p=3/7$ and are strictly below one on their
stated intervals. No claim of attainability is made. The comparison
argument uses positivity of the success map, while the bound for its
output on $\omega$ uses the LOCC Schmidt-number constraint.

Grouping copies into pairs also gives, for every $n\ge2$,
\[
 \operatorname{SN}(\omega^{\otimes n})
        \le 7^{\lfloor n/2\rfloor}3^{n\bmod2}<3^n,
 \qquad G_{3^n}(\omega^{\otimes n})=0,
 \qquad G_{3^n}^{\rm a}(\omega^{\otimes n})\ge(3/7)^n.
\]
In particular,
$E_c^{(0,\infty)}(\omega)\le\tfrac12\log_2 7<\log_2 3$.
An arbitrary collective LOCC success branch into a $3^n\times3^n$ space
has FEF at most $(7/9)^{\lfloor n/2\rfloor}$, since it cannot increase
the displayed Schmidt number. This last statement concerns output at the
full dimension $3^n$; it does not exclude distillation to a smaller
maximally entangled state.

\subsection{Minimality and prior-art boundary}

The failure just proved first becomes possible in square local dimension
three among states with positive equal $G$ roofs. By Gour and Guo's
classification, a two-qubit state with $G_2=G_2^{\rm a}>0$ has coefficient
support $X(\mathbb CI+\mathcal N)$, where $X$ is invertible and
$\mathcal N$ is a nilpotent linear space of dimension at most one.
A single nilpotent matrix is triangularizable. Thus invertible local
filtering carries its support into a common triangular flag, and the
universal tensor laws proved above apply. Dimension one is trivial.
Our explicit dimension-three example proves that this minimality is sharp.

Gour and Guo's Theorem 7 in
\href{https://quantum-journal.org/papers/q-2018-08-13-81/}{Quantum 2, 81 (2018)}
already characterizes positive equality of the two $G$ roofs using
nilpotent linear spaces. That single-state result is credited here.
The nilpotent pencil and the distinction between scalar and matrix-valued
pencils also have prior matrix-theory antecedents in Christensen's paper.
Terhal and Horodecki,
\href{https://arxiv.org/abs/quant-ph/9911117}{arXiv:quant-ph/9911117},
already established that mixed-state Schmidt number need not multiply
under tensor products. Yue and Chitambar's subsequent zero-error-cost
work gives strong forms of that nonmultiplicativity. The additions here
are the universal factorization laws forced by the spectral-saturator
support, and the exact positive-flat-roof counterexample separating a
common triangular flag from a general nilpotent linear space, together
with its quantitative filtering obstruction. Targeted primary-source
searches did not locate these precise statements; they do not certify
worldwide priority.

These tensor and convex-roof theorems are written exact proofs. The
companion verification script checks determinant identities, ensemble
normalization, actual marginals, ranks, and random ensemble identities;
those checks do not by themselves formalize Schmidt-number convex roofs
or LOCC in Lean. Formal coverage is reported separately.
